\chapter{Glossitis}

\section*{Definition}
Glossitis is an inflammatory condition of the tongue characterized by swelling, redness, loss of filiform papillae (atrophic glossitis), and often pain or burning; acute glossitis involves rapid onset of edema, while chronic glossitis presents as smooth, red, depapillated mucosa.

\section*{Pathophysiology}
Inflammatory infiltration of the tongue mucosa leads to atrophy of the filiform papillae, thinning of the epithelium, and vascular congestion. Nutritional deficiency states (particularly B vitamins) impair epithelial cell turnover, causing loss of the normal papillae architecture and exposure of the underlying erythematous mucosa.

\section*{Causes}
\begin{itemize}
  \item \textbf{Nutritional deficiencies:} Vitamin B12 (pernicious anemia), folate, iron (Plummer-Vinson syndrome), niacin (pellagra), riboflavin, pyridoxine
  \item \textbf{Infectious:} Candidiasis (acute atrophic candidiasis), bacterial glossitis (Streptococcus, Actinomyces), HSV stomatitis
  \item \textbf{Allergic/irritant:} Contact allergy (toothpastes, mouthwashes, foods, dental materials), trauma (burn, biting)
  \item \textbf{Systemic disease:} Sjogren syndrome, celiac disease, Crohn disease, Behcet disease, syphilis (luetic glossitis)
  \item \textbf{Drug-induced:} Chemotherapy (mucositis), antibiotics, ACE inhibitors, NSAIDs
\end{itemize}

\section*{When to See a Doctor}
See your doctor if:
\begin{itemize}
  \item Persistent tongue burning, pain, or swelling lasting > 2 weeks
  \item Smooth, red, "beefy" tongue appearance (suspect nutritional deficiency)
  \item Dysphagia, odynophagia, or hoarseness
  \item Systemic symptoms: fatigue, pallor, paresthesias, or weight loss
  \item Acute tongue swelling with respiratory compromise (angioedema -- emergency)
\end{itemize}

\section*{Related Symptoms}
\begin{itemize}
  \item Tongue pain and burning (glossodynia)
  \item Erythema and edema of the tongue
  \item Smooth, atrophic, depapillated dorsal surface
  \item Dysgeusia (altered or metallic taste)
  \item Xerostomia (if associated with Sjogren or radiation)
\end{itemize}
\textbf{Important:} Atrophic glossitis ("bald tongue") is a hallmark of vitamin B12 deficiency and may precede the development of anemia and neurologic symptoms by months.

\section*{Self-Care / Home Management}
\begin{itemize}
  \item Avoid spicy, acidic, or excessively hot foods that exacerbate irritation
  \item Soft, bland diet during acute inflammation
  \item Avoid alcohol, tobacco, and alcohol-based mouthwashes
  \item Oral hygiene with a soft-bristled toothbrush
  \item Saline mouth rinses (1/2 tsp salt in 8 oz water) 2--3 times daily
  \item Over-the-counter multivitamin with B-complex, folate, and iron if deficiency suspected
\end{itemize}

\section*{How Your Doctor Will Evaluate You}
\begin{itemize}
  \item Inspection of tongue: color, texture, papillae status, swelling, and ulceration
  \item Complete blood count with RBC indices, serum B12, folate, ferritin, iron, and TIBC
  \item KOH preparation and fungal culture if candidiasis suspected
  \item Patch testing for contact allergy if history suggests allergen exposure
  \item Anti-parietal cell and intrinsic factor antibodies for pernicious anemia
  \item Biopsy for non-healing or asymmetric lesions to exclude carcinoma
  \item Serologic testing for celiac disease (tTG-IgA) if associated GI symptoms
\end{itemize}

\section*{Treatment Options}
\begin{itemize}
  \item Nutritional replacement: IM vitamin B12 (pernicious anemia), oral folate 1 mg daily, oral iron (ferrous sulfate 325 mg daily)
  \item Topical corticosteroids (triamcinolone dental paste 0.1\%) for idiopathic or allergic glossitis
  \item Antifungal therapy (nystatin suspension, fluconazole) for candidal glossitis
  \item Avoidance of identified allergens or irritants
  \item Management of underlying systemic disease (gluten-free diet for celiac, immunosuppression for Crohn)
  \item Sialogogues (pilocarpine, cevimeline) for xerostomia-associated glossitis
  \item Pain management: topical viscous lidocaine or systemic analgesics
\end{itemize}

\clearpage
